\section{Discussion}

Figures~\ref{fig:normal_damp},~\ref{fig:crit_damp}, and~\ref{fig:osc_damp} show the voltage signals of the components
agree with the theoretical predictions.
The critically damped system discharges quicker than the normally damped configuration, and the oscillatory system
shows the electric potential transferring between the magnetic and electric fields.
Furthermore, the theoretical models for $V_R$ and $V_C$ align almost exactly to the recorded experimental signals.
Small differences between the experimental and theoretical results here can be explained by the resolution of the
evaluated time step, $\Delta t$.
As $\Delta t$ converges to zero, the difference between the theory and experiment would decrease.
The $V_R$ and $V_C$ signals attenuate much faster in the critically damped case.
This is due to the current decreasing faster as a result of the reduced resistance in the circuit.

However, in examining the energetics in Figures~\ref{fig:case2_energy} and~\ref{fig:case3_energy}, the conserved energy
unexpectedly rose above the initial total energy.
This suggests the sum of cumulative heat dissipation energy and total energy did not agree with theoretical predictions.
Since no external energy sources existed in the circuit, the increase in energy can be explained by a few
contributing factors.
In the oscillatory case, the energy transfer between electric and magnetic fields decreased in magnitude over time but
never reached
zero during the cycle.
In addition, the accumulated heat dissipation, $E_{HT}$, increased rapidly and unexpectedly rose higher than the
initial energy of the system.
This can be explained by the fact that the whole cycle was not captured from initial discharge, but indeed
began slightly after due to the software's trigger settings.


$E_T$ and $E_{HT}$ show to have correlated signals in both cases.
When $E_T$ drops, $E_{HT}$ rises, and when $E_T$ begins to plateau, so does $E_{HT}$.
This is because the total energy of the components transferred into the current through the system, thus resulting
in joule heating.
In Case 2, $E_{HT}$ plateaus quickly because current rapidly attenuates from critically damped discharge.
In Case 3, $E_{HT}$ gradually plateaus because the induced current from oscillations takes longer to attenuate.

Also, though a small contribution, the model does not take into account the thermal heat dissipated through the wires
connecting the circuit, only that of the system impedance of the individual components.
As temperatures change, so do the internal resistances of the components, further affecting the
calculated result due to changing current flow.

Analysing the frequency response of the oscillatory case, we found that the experimentally determined frequency was
consistent with the theoretical frequency taking into account the system resistance, $R$.
The theoretical natural frequency of the system was found to be slightly higher and just outside
consistency margins.
This suggests the damped frequency is slightly lower than the natural frequency of the system due to the
contributing damping resistance, $R$.


\section{Conclusion}

From our experimental analysis, the transient response of the LRC circuit under normal, critical, and oscillatory
damped configurations demonstrated the relationship between electric and magnetic fields during a discharge event.
The electric potential clearly showed that energy was exchanged back and forth between the components, and the
experimental
signals strongly agreed with the theoretical prediction curves.
In addition, the frequency calculations were consistent with experimental observations.

The energetics plots
demonstrated an unexpected increase in conserved energy, but this was explained by the data being captured too far
ahead
of the initial discharge event.
This erroneous reduction in the initial total energy stored in the
circuit led to a pseudo upward trend in the conserved energy.
By adjusting the trigger settings to be more sensitive to the falling edge of the signal, a full capture of the
discharge cycle, from fully
charged to
fully discharged, would
have shown a higher initial total
energy, and thus shown the conservative energy to gradually decrease.

